\section{Experiments}
\label{sec:experiments}

The experimental section validates the claims made in the introduction.

\subsection{Datasets and Baselines}
Specify the datasets used for evaluation. Detail the baseline models for fair comparison.

\subsection{Results and Analysis}

Results are typically summarized in tables. Table \ref{tab:results} demonstrates a standard formatting approach using \texttt{booktabs} for clean, professional rules without vertical lines.

\begin{table}[H]
    \centering
    \caption{Performance comparison on standard benchmarks.}
    \label{tab:results}
    \begin{tabular}{lccc}
        \toprule
        \textbf{Method} & \textbf{Precision} & \textbf{Recall} & \textbf{F1-Score} \\
        \midrule
        Baseline A      & 0.782              & 0.810           & 0.795             \\
        Baseline B      & 0.824              & 0.801           & 0.812             \\
        \textbf{Ours}   & \textbf{0.891}     & \textbf{0.885}  & \textbf{0.888}    \\
        \bottomrule
    \end{tabular}
\end{table}

% \begin{figure}[t]
%     \centering
%     \includegraphics[width=\linewidth]{figures/arch.pdf}
%     \caption{Overview of the proposed system architecture.}
%     \label{fig:architecture}
% \end{figure}

Figures should be referenced in the text. For instance, see Figure \ref{fig:architecture} for the system overview (commented out in the source code).
